\chapter{Язык между модулем и дроном}\label{ch:uart}

Модуль и автопилот соединены обычным последовательным кабелем: линии TX, RX, питание и земля. Скорость 115\,200 бод, на стороне дрона это устройство \texttt{/dev/ttyS0}. Поверх кабеля идёт собственный простой протокол без названия. Мы восстановили его с трёх сторон: по коду модуля, по коду автопилота и по живому обмену. \pr{}

\section{Кадр модуля}\label{sec:frame}

Модуль отправляет дрону кадр длиной \textbf{22 байта}. Первый байт всегда \texttt{CD}, последний --- контрольная сумма (сумма байтов 1--20 по модулю 256).

\begin{figure}[tbp]
\centering
\begin{tikzpicture}[x=7.15mm, y=1mm,
  c/.style={draw=rule, minimum width=7.15mm, minimum height=7mm, font=\ttfamily\scriptsize, inner sep=0pt},
  lab/.style={font=\sffamily\tiny, color=inkgray, align=center}
]
\foreach \i/\t/\col in {0/CD/accentsoft, 1/02/warnfill, 2/00/stopfill, 3/00/stopfill, 4/ff/accentsoft, 5/fc/accentsoft, 6/00/warnfill, 7/16/warnfill, 8/01/accentsoft, 9/28/accentsoft, 10/0d/warnfill, 11/f3/warnfill, 12/01/stopfill, 13/00/paper, 14/00/paper, 15/00/paper, 16/00/paper, 17/00/paper, 18/1f/accentsoft, 19/01/accentsoft, 20/00/paper, 21/8f/warnfill}
  \node[c, fill=\col] at (\i+0.5, 0) {\t};
\foreach \i in {0,...,21} \node[lab] at (\i+0.5, -7) {\i};
\draw[decorate, decoration={brace, amplitude=2pt, mirror}] (0.05,-11) -- (0.95,-11);
\node[lab] at (0.5,-17) {заголовок};
\draw[decorate, decoration={brace, amplitude=2pt}] (1.05,5) -- (1.95,5);
\node[lab] at (1.5,11) {тип};
\draw[decorate, decoration={brace, amplitude=2pt}] (2.05,5) -- (3.95,5);
\node[lab] at (3,11) {номер\\метки};
\draw[decorate, decoration={brace, amplitude=2pt}] (4.05,5) -- (5.95,5);
\node[lab] at (5,11) {$x$};
\draw[decorate, decoration={brace, amplitude=2pt}] (6.05,5) -- (7.95,5);
\node[lab] at (7,11) {$y$};
\draw[decorate, decoration={brace, amplitude=2pt}] (8.05,5) -- (9.95,5);
\node[lab] at (9,11) {дистан-\\ция};
\draw[decorate, decoration={brace, amplitude=2pt}] (10.05,5) -- (11.95,5);
\node[lab] at (11,11) {угол};
\draw[decorate, decoration={brace, amplitude=2pt}] (12.05,5) -- (12.95,5);
\node[lab] at (12.5,11) {флаг};
\draw[decorate, decoration={brace, amplitude=2pt}] (18.05,5) -- (19.95,5);
\node[lab] at (19,11) {ToF,\\мм};
\draw[decorate, decoration={brace, amplitude=2pt, mirror}] (21.05,-11) -- (21.95,-11);
\node[lab] at (21.5,-17) {сумма};
\end{tikzpicture}
\caption{Живой кадр модуля, метка 0 под камерой. Нулевые байты 13--17 и 20 зарезервированы.}\label{fig:frame}
\end{figure}

\begin{tabularx}{\textwidth}{@{}c p{3.1cm} Y@{}}
\toprule
\hdr{Байты} & \hdr{Поле} & \hdr{Смысл} \\ \midrule
0 & заголовок & Всегда \texttt{CD}. \\
1 & тип & \texttt{02} --- кадр о метке. \\
2--3 & номер метки & 16 бит, старший байт первым. Метка 0 дала \texttt{00 00}, метка 13 дала \texttt{00 0d}. \\
4--5, 6--7 & смещения $x$, $y$ & Знаковые 16 бит, насколько камера сдвинута от центра метки. \\
8--9 & расстояние & До метки, знаковые 16 бит. \\
10--11 & угол & Поворот относительно метки, в десятых долях градуса. \\
12 & флаг & 1 --- метка найдена, 0 --- нет. \\
13--14 & счётчик & 16 бит (в наших кадрах 0). \\
15--17, 20 & резерв & Нули. \\
18--19 & ToF & Расстояние до пола в миллиметрах, \textbf{младший байт первым}. \\
21 & сумма & Сумма байтов 1--20 по модулю 256. \\ \bottomrule
\end{tabularx}

\vspace{4pt}
\pr{} Кадр снят дважды: сначала в памяти работающей программы модуля, потом уже \emph{внутри} автопилота на дроне. Данные совпали байт в байт.

\begin{zametka}[Проверка в цифрах]
Метка 0 под камерой, наблюдаемый кадр: расстояние 296, ToF 287 мм (почти одно число), угол 3571 (то есть 357,1\textdegree, почти ноль, модуль был повёрнут ровно), смещения $-4$ и $22$ (камера почти над центром метки). Всё согласуется. \pr{}
\end{zametka}

\section{Что делает дрон с кадром}

В автопилоте работает конечный автомат приёма. Он принимает байты из \texttt{/dev/ttyS0}, узнаёт заголовок, накапливает данные и проверяет сумму. Правильные данные копируются в буфер, и тип 2 запускает разбор метки. В памяти работающего \texttt{pilot\_f} мы прочитали этот буфер: в нём лежали те же 20 байт данных, что отправлял модуль. У драйвера UART счётчик принятых байт рос: «20\,152, 20\,262». \pr{}

\section{Версия заголовка: важная тонкость}

Приёмник в старой прошивке дрона (март 2023) и приёмник в образе 1.2.8 (2025) \emph{различаются}.

\begin{tabularx}{\textwidth}{@{}l Y@{}}
\toprule
\hdr{Прошивка автопилота} & \hdr{Заголовок кадра от модуля} \\ \midrule
\texttt{pilot\_f}, март 2023 (на нашем дроне) & Один байт: \texttt{CC} (14 байт данных) или \texttt{CD} (20 байт данных). \\
\texttt{pilot}, образ 1.2.8 (2025) & Два байта: \texttt{CC CC}, \texttt{CD CD} или \texttt{CE CE}. \\ \bottomrule
\end{tabularx}

\vspace{4pt}
Модуль с программой, которая шлёт один байт \texttt{CD}, дружит со старой версией автопилота. Для образа 1.2.8 нужна программа модуля с двумя байтами заголовка. \pr{} (оба приёмника прочитаны); живая проверка сделана только со старой версией.

\begin{vnimanie}[Если делаете свою камеру]
Формат кадра нужно подбирать под \emph{версию автопилота на дроне}. Сначала узнайте её (глава~\ref{ch:kak}), потом пишите кадр.
\end{vnimanie}

\section{Команды дрона к модулю}

Дрон тоже шлёт модулю кадры. Они начинаются с двух байт \texttt{AD FA}, затем идёт команда, данные и сумма по данным. Разбор сделан по коду модуля (\texttt{ckh831\_protocol\_recv}).

\begin{tabularx}{\textwidth}{@{}c Y@{}}
\toprule
\hdr{Команда} & \hdr{Что делает модуль} \\ \midrule
1 & Выбирает режим (0 --- метки, 1 --- блобы, 2 --- цвет). \\
2 & Включает выход GPIO (линии \texttt{0xE2}, \texttt{0xE3}). \\
3 & Задаёт цвет светодиодов WS2812. \\
4, 6 & Устанавливает ШИМ на каналах 1 и 2 (во второй команде значение 16-битное). \\
5 & Настраивает цвет для режима «цвет», запись в JSON. \\
9 & Отвечает двумя байтами \texttt{F9 66} (проверка связи). \\
10 & Делает снимок кадра. \\
\texttt{F1} & Устанавливает или сбрасывает внутренний флаг. \\
\texttt{FA} & Пересылает полученные байты на порт UART0 модуля (\texttt{ttyS0}, разъём U0TX/U0RX). Через неё идут команды внешнему зуммеру (глава~\ref{ch:pult}). \pr{} (проверено на слух) \\ \bottomrule
\end{tabularx}

\vspace{4pt}
\vyv{} Это разбор по коду. Какие из команд дрон шлёт на деле, мы не снимали: за пять минут на дроне в сторону модуля ушло 360 байт. Команды 3 и \texttt{FA} мы потом посылали на живой модуль и увидели результат (светодиод, зуммер). \pr{}

\section{Что не выяснено}

Другие типы кадров (1, 4, 5, \texttt{F9}) и кадр \texttt{CE CE} в автопилоте; смысл счётчика в байтах 13--14. Зачем модуль пересылает байты на второй порт, теперь ясно: это канал к внешним устройствам вроде зуммера (глава~\ref{ch:pult}); кто пользуется им на борту штатно, не знаем. \nei{}
